\section{Retrieval-Augmented Generation}
\label{sec:bg-rag}

Retrieval-augmented generation supplies a model with material fetched at
inference time rather than relying on pre-training
alone~\cite{lewis2020rag}. A retriever selects passages from a corpus and the
generator conditions on them together with the query. For software the appeal
is immediate: a repository at a particular commit is knowledge no pre-trained
model can have.

Modern retrieval is typically dense. Karpukhin et al.\ showed that a dual
encoder mapping questions and passages into a shared embedding space
outperforms sparse baselines on open-domain question
answering~\cite{karpukhin2020dpr}. Three parameters recur throughout
Chapter~\ref{chap:approach}: how the corpus is chunked, how the query is
formed, and how many chunks are kept.

The property that matters for code review is what similarity does and does not
guarantee. Similarity in embedding space is a claim about resemblance, not
about consequence. A function that resembles the changed code will rank highly
whether or not the change affects it; the caller that will actually break may
share no surface features with the callee and rank poorly. Retrieval optimises
for looking related; review needs what \emph{is} related.

Graph-structured retrieval is one response.
GraphRAG~\cite{edge2024graphrag} builds an entity graph over the corpus to
answer questions about how documents connect rather than about any single
document. RepoCoder~\cite{zhang2023repocoder} alternates retrieval and
generation for repository-level completion, on the finding that a single
similarity query over a diff retrieves poorly. Both are admissions about the
limits of similarity as a relevance criterion and both point at the comparison
this thesis makes: if similarity retrieval were adequate, neither iterative
re-querying nor graph construction would be necessary.

The retrieval arm in this thesis is a deliberately conventional instance of the
paradigm (chunked corpus, cosine similarity, fixed top-$k$), included not
as a contribution but as the comparison that makes the graph arm
interpretable.
